
\paragraph{Helmholtz–Hodge Stress.}
Recall the shear stress field of Problem 2 takes the form:
\[
\bm{\tau}=G\Big[\ c_{\vartheta}\big(\nabla \WarpTwistIesan+\FrameBasis\times\bm{r}\big)
+\nabla (\bm{\WarpShearIesan}\cdot\bm{b}_{\Section})
-\nu\bigl(\operatorname{dev}(\bm{r}\otimes\bm{r})-\bm{r}\otimes\CentroidLocation \bigr)\,\bm{b}_{\Section}
\Big] \\
\]
Introduce the solenoidal b-\emph{Poisson tensor}
\[
\bm{\Pi} = \operatorname{dev}\bm{r}\otimes\bm{r} - \bm{r}\otimes\CentroidLocation
\]
Define the Poisson vector field
\[
\begin{aligned}
\bm{\pi}
&=\bm{\Pi}\,\bm{b}_{\Section}
\end{aligned}
\qquad\text{so}\qquad
\begin{aligned}
\operatorname{div} \bm{\pi} &= \,2\,\big(\bm{r}-\CentroidLocation\big) \cdot \bm{b}_{\Section} \\
\operatorname{curl}\bm{\pi}
&= 2\,(\FrameBasis\times\bm{r})\cdot\bm{b}_{\Section}.
\end{aligned}
\]
$$
\operatorname{grad} \bm{\pi} = \bm{b}_{\Section} \otimes \bm{r} + \big( (\bm{r} - \CentroidLocation) \cdot \bm{b}_{\Section} \big) \mathbf{1} - \frac{2}{3} \bm{r}\otimes \bm{b}_{\Section}
$$
%
The traction–free boundary condition for \eqref{eq:bvp-warp-shear-iesan} is compactly
\(
\partial_{\bm\nu}\WarpShearIesan = \nu\,\bm{\pi}\cdot\bm{\nu}.
\)
%
In terms of \(\bm{\pi}\) the stress is:
\[
\bm{\tau}=G\big(\nabla\WarpShear-\nu\,\bm{\pi}\big)
\]
Decompose \(\bm{\pi}\) via a scalar potential \(\WarpShearPotentialPoisson\), a stream function \(\zeta\), and a harmonic field \(\bm{h}\):
\[
\bm{\pi}=\nabla\WarpShearPotentialPoisson + \FrameBasis\times\nabla\zeta + \bm{h},
\]
where \(\bm{h}\) is harmonic Neumann (\(\operatorname{div}\bm{h}=\operatorname{curl}\bm{h}=0,\ \bm{h}\cdot\bm\nu=0\) on \(\partial\Section\)) and \(\WarpShearPotentialPoisson\) solves the Neumann problem
\[
\left\{
\begin{aligned}
\Delta\WarpShearPotentialPoisson&=\operatorname{div}_{\Section}\bm{\pi} , % =2(\bm{r}-\CentroidLocation)\cdot\bm{b}_{\Section},
\\
\partial_{\bm\nu}\WarpShearPotentialPoisson&=\bm{\pi}\cdot\bm\nu
&\quad\text{on } \SectionBoundary,
\end{aligned}
\right.
\]
with the usual zero–mean gauge. The stream \(\zeta\) solves the Dirichlet problem:
\[
\begin{cases}
\Delta \,\zeta \;=\; \operatorname{curl}\bm{\pi} & \text{in }\Section,\\[2pt]
\zeta\big|_{\SectionBoundary_k} \;=\; 0 & \text{on }\Gamma_k,\;k\in \{1,...,\dim\bm{h}\}.
\end{cases}
\]

Define the Helmholtz potentials for \(\bm{\tau}\) by
\[
\WarpShearPotential:=G\Big(c_{\vartheta}\,\WarpTwistIesan
\;+\;\WarpShearIesan
\;-\;\nu\,\WarpShearPotentialPoisson\Big),\qquad
\WarpShearStream:=-\,G\Big(\nu\,\zeta+\tfrac12\,r^2\,c_{\vartheta}\Big).
\]
Then the shear admits the Helmholtz–Hodge decomposition \citep{cantarella2002vector}
\[
\boxed{
\ \bm{\tau}=\nabla \WarpShearPotential
\;+\; \FrameBasis\times\nabla\WarpShearStream
\;-\;G\nu\,\bm{h}\ }.
\]

\noindent\emph{Consequences.}
Using \(\Delta\WarpShearIesan=-2(\bm{b}_{\Section}\cdot\bm{r}+b_\varepsilon)\) and \(\Delta \WarpShearIesan = 0\)
\[
\operatorname{div}\bm{\tau}
=\Delta\WarpShearPotential
=G\Big(\Delta(\bm{\WarpShearIesan}\cdot\bm{b}_{\Section})
-\nu\,\Delta\WarpShearPotentialPoisson\Big)
=-E\,(\bm{b}_{\Section}\cdot\bm{r}+b_\varepsilon),
\]
\[
\operatorname{curl}\bm{\tau}
=\operatorname{curl}\big(\FrameBasis\times\nabla\WarpShearStream\big)
=-\,G\nu\,\Delta\zeta\;-\;2G\,c_{\vartheta}
=\,-\,G\nu\,\operatorname{curl}\bm{\pi}\;-\;2G\,c_{\vartheta}.
\]
Note that \(\operatorname{curl}(\nabla \WarpShearPotentialPoisson)=\operatorname{curl} \bm{h}=0\). On \(\SectionBoundary\),
\[
\partial_{\bm\nu}\WarpShearPotential
+(\nabla\WarpShearStream)\cdot(\FrameBasis\times\bm\nu)
-G\nu\,\bm{h}\cdot\bm\nu
=0,
\]
which is equivalent to the Saint–Venant free–side condition for \(\WarpShearIesan\) together with \eqref{eq:bvp-warp-twist-iesan}.
